% --------------------------------------------------------
% AI USAGE
% --------------------------------------------------------
\noindent Throughout the project, the team incorporated AI across multiple stages to assist with various tasks. Importantly, AI was never used as a replacement for human effort or core decision-making. Instead, it served exclusively as a supplementary tool to streamline team workflows and optimize overall time management. 

\subsection*{Some of the examples using AI in this project:}
\begin{itemize}
    \item \textbf{Architectural Brainstorming:} In the early planning stages, AI was leveraged to help generate ideas for potential design patterns and library integrations. Team then carefully weighed these options against established non-functional requirements before finalizing the architecture. 
    \item \textbf{Debugging Assistance:} Dealing with undocumented Avalonia UI errors or sudden backend disconnects sometimes meant untangling obscure error codes. AI was used as a supplementary tool to quickly parse these codes, though diagnosing the root cause within our specific architecture and developing the actual fixes remained a manual team effort. 
    \item \textbf{Formatting and Documentation:} To streamline team workflow, it was decided to outsource repetitive formatting tasks to AI tools. This included generating UML diagram sketches from our existing C\# codebase to then format those in a specific editor and resolving some of tricky LaTeX compilation errors in the report, which saved valuable time without affecting the design or logic of the core system. 
    \item \textbf{UI/UX Inspiration:} Various stylistic directions were explored for project logos and interface by generating mood boards with AI image tools. Even though the final application relies exclusively on original assets, these early visualizations helped steer overall design choices. 
    \item \textbf{Code Review and Refactoring:} During the refactoring phase, AI acted as an extra set of eyes to catch missed references. When renaming methods within optimization modules, for instance, it helped verify that all old instances were properly updated, effectively reducing minor human errors without altering the underlying code logic.
\end{itemize}
